% =====================================================================
%  书末：后记（あとがき，原书 p.373–374）
%  “単行本化にあたっての追記”（原书 p.375–384）另见 chunks/afterword.tex
% =====================================================================
\chapter*{后记}
\addcontentsline{toc}{chapter}{后记（原书 p.373）}
\markboth{后记}{后记}

% [原文 p.373]
正如开头所述，辛几何的全貌才刚刚开始显露。本书不过是现时点已出现的部分中、笔者所见到的部分中的又一部分，还有许多未能言及。

与动力系统研究的联系未能充分叙述，尤其遗憾。

本书着力点之一的量子化，笔者也还在学习之中，知识有所偏重，这也反映在本书中。

奇点论、微局部分析、表示论、可积系统等，各自都是极重要的话题，与辛几何的关系也很深，但本书未能多作叙述。

关于辛拓扑，近来在各种场合已有许多论述，因此与它的重要性相比，本书对它的处理较轻。

在正文中设置多到几乎过分的“注”，一面引用文献、一面尽量触及更多的领域，正是因为想略微弥补上述缺陷。

辛几何虽未完成，却已成长为巨大的领域，一人叙述其全貌是不可能的。许多人按照各自的想法去讲述、研究各个侧面，才是重要的——我认为如此；因此，尽管作为讲座的一册具有被要求正统性的读物性质，笔者仍不掩饰自己的喜好与偏见。等待读者的批评，就此搁笔。
